\documentclass[11pt]{article}
\usepackage{amsmath,amssymb,amsthm,mathrsfs}
\usepackage[margin=1.1in]{geometry}
\usepackage[hidelinks]{hyperref}

\theoremstyle{plain}
\newtheorem{theorem}{Theorem}
\newtheorem{lemma}[theorem]{Lemma}
\newtheorem{proposition}[theorem]{Proposition}
\newtheorem{corollary}[theorem]{Corollary}
\theoremstyle{definition}
\newtheorem{definition}[theorem]{Definition}
\newtheorem{hypothesis}[theorem]{Hypothesis}
\theoremstyle{remark}
\newtheorem{remark}[theorem]{Remark}

\newcommand{\R}{\mathbb{R}}
\newcommand{\C}{\mathbb{C}}
\newcommand{\Z}{\mathbb{Z}}
\newcommand{\Lam}{\Lambda}
\DeclareMathOperator{\spec}{spec}
\DeclareMathOperator{\rank}{rank}
\DeclareMathOperator{\inertia}{inertia}
\DeclareMathOperator{\ran}{ran}

\title{The Gram double pencil: a spectral census for the zeros of $\zeta$,\\
an Euler anchor, and a bank-generated warp (Version 2)}
\author{}
\date{2026-08-14}

\begin{document}
\maketitle

\begin{abstract}
We attach to each ordinate window a pair of Hankel matrices computed by
contour integration of $\xi'/\xi$, and prove that the number of zeros off the critical
line in the window equals the negative inertia of the first matrix (Theorem
\ref{thm:inertia}); the remaining spectral data of the pair counts the zeros with
multiplicity, their distinct support, and their ordinates (Theorem \ref{thm:census}).
Both statements are unconditional. On symmetric pairs of windows, the location statement
is equivalent to a Stieltjes-type positivity condition in the quotient chart $w=z^2$
(Proposition \ref{prop:stieltjes}). We then construct,
in the half plane of absolute convergence and with no reference to zeros, an
Euler-anchored positive-definite double pencil (Proposition \ref{prop:anchor}), an
explicit finite-dimensional regulator (Theorem \ref{thm:regulator}), and an exact
bank-generated warp for every regular Hermitian-definite pencil path (Theorems
\ref{thm:warp} and \ref{thm:pullback}). These finite-dimensional results do not identify
the terminal pencil with the contour-moment pencil. That endpoint identification is
isolated as the seating hypothesis in \S\ref{sec:seat}; the analytic defect calculation
needed to test it is stated in \S\ref{sec:obligations}.
\end{abstract}

\section{Setup}

Let $\xi$ be the completed Riemann zeta function and put
\[
  A(z) := \xi\!\left(\tfrac12 + iz\right),
\]
an entire function of order one, even ($A(-z)=A(z)$, the functional equation), real on
$\R$, whose zeros are
\[
  z_\rho=-i\!\left(\rho-\tfrac12\right)
  =\gamma_\rho+i\!\left(\tfrac12-\beta_\rho\right)
\]
for
$\rho = \beta_\rho + i\gamma_\rho$ a nontrivial zero of $\zeta$; in particular
\[
  \rho \text{ on the critical line} \iff z_\rho \in \R .
\]
Nonreal zeros of $A$ occur in conjugate pairs $z,\bar z$ with equal multiplicities
(conjugate symmetry of $\xi$ on the real axis).

\begin{definition}[window moments]\label{def:moments}
Let $W=(a,b)$ with $A$ nonvanishing on $\partial W \times \{|y|\le Y\}$, $Y>\tfrac12$,
and let $\partial R_W$ be the positively oriented rectangle boundary. Put
\[
  \mu_k(W) := \frac{1}{2\pi i}\oint_{\partial R_W} z^k\,\frac{A'(z)}{A(z)}\,dz,
  \qquad k \ge 0,
\]
and for $n\ge 1$,
\[
  H_n(W) := \bigl(\mu_{\alpha+\beta}(W)\bigr)_{\alpha,\beta=0}^{n-1},
  \qquad
  H_n^{(1)}(W) := \bigl(\mu_{\alpha+\beta+1}(W)\bigr)_{\alpha,\beta=0}^{n-1}.
\]
\end{definition}

By the argument principle, with $x_1,\dots,x_r$ the distinct real zeros of $A$ in the
window (multiplicities $a_i\ge 1$) and $z_j,\bar z_j$, $j=1,\dots,q$, the distinct
conjugate pairs (equal multiplicities $b_j \ge 1$),
\begin{equation}\label{eq:moments}
  \mu_k(W) \;=\; \sum_{i=1}^r a_i x_i^k \;+\; \sum_{j=1}^q b_j\bigl(z_j^k + \bar z_j^{\,k}\bigr).
\end{equation}
All $\mu_k(W)$ are real and computable from $A$ alone; no knowledge of the zeros is used.
Write $m := r + 2q$ for the number of distinct support points.

\section{The inertia identity}

\begin{theorem}[inertia]\label{thm:inertia}
For every $n \ge m$,
\[
  \inertia H_n(W) \;=\; \bigl(r+q,\; q,\; n-r-2q\bigr),
\]
where $\inertia$ denotes (positive, negative, zero) eigenvalue counts. In particular
\[
  n_-\bigl(H_n(W)\bigr) \;=\; q,
  \qquad\text{and}\qquad
  H_n(W) \succeq 0 \iff q = 0,
\]
i.e.\ the negative inertia counts exactly the off-critical-line conjugate pairs.
\end{theorem}

\begin{proof}
Let $v(t) := (1,t,\dots,t^{n-1})^{T}$. By \eqref{eq:moments},
\[
  H_n \;=\; \sum_{i=1}^r a_i\, v(x_i)v(x_i)^{T}
  \;+\; \sum_{j=1}^q b_j\Bigl[v(z_j)v(z_j)^{T} + v(\bar z_j)v(\bar z_j)^{T}\Bigr].
\]
Write $v(z_j) = A_j + iB_j$ with $A_j,B_j\in\R^n$; then $v(\bar z_j)=A_j-iB_j$ and
\begin{equation}\label{eq:minus}
  v(z_j)v(z_j)^{T} + v(\bar z_j)v(\bar z_j)^{T}
  \;=\; 2A_jA_j^{T} - 2B_jB_j^{T},
\end{equation}
the minus sign being the whole mechanism. Hence $H_n = UJU^{T}$ with
\[
  U = \bigl[v(x_1)\cdots v(x_r)\;A_1\,B_1\cdots A_q\,B_q\bigr],
  \qquad
  J = \operatorname{diag}(a_1,\dots,a_r,\,2b_1,-2b_1,\dots,2b_q,-2b_q),
\]
so $\inertia J = (r+q,q,0)$.

$U$ has full column rank $m$: the complex matrix
$V=[v(x_1),\dots,v(x_r),v(z_1),v(\bar z_1),\dots]$ is Vandermonde on $m$ distinct nodes
with $n\ge m$, so its leading $m\times m$ minor equals
$\prod_{\alpha<\beta}(s_\beta-s_\alpha)\ne 0$ and $\rank V=m$; replacing each pair
$(v(z_j),v(\bar z_j))$ by $(A_j,B_j) = \bigl(\tfrac{v(z_j)+v(\bar z_j)}{2},
\tfrac{v(z_j)-v(\bar z_j)}{2i}\bigr)$ is an invertible change of columns, so
$\rank U = m$.

Take a thin $QR$ factorisation $U=QR$ with $Q^{T}Q=I_m$ and $R$ invertible, and extend
$Q$ to an orthogonal $\widetilde Q = [Q,Q_\perp]$. Then
\[
  \widetilde Q^{T} H_n \widetilde Q =
  \begin{pmatrix} RJR^{T} & 0\\ 0 & 0_{n-m}\end{pmatrix},
\]
and by Sylvester's law of inertia $\inertia(RJR^{T}) = \inertia(J)$. The zero block
contributes $n-m$ null directions.
\end{proof}

\begin{remark}
The hypothesis $n \ge m$ is necessary: an off-line pair is invisible to a Hankel matrix
smaller than the distinct support count. Since $\mu_0(W)$ is the number of zeros in $W$
with multiplicity and $m \le \mu_0(W)$, the choice
$n := \max\{1,\mu_0(W)\}$ is admissible and computable from the contour data alone.
\end{remark}

\begin{remark}[collisions]
If a conjugate pair approaches $\R$, the corresponding negative eigenvalue tends to $0^-$
and at collision becomes a null direction; the invariant through such an event is the
inertia, not any individually ordered eigenvalue. Consequently the correct target is
non-negativity $H_n \succeq 0$ (equivalently $n_-=0$), never strict positivity: the
$n-m$ structural nulls and the collision nulls are expected.
\end{remark}

\section{The census, and the equivalent form of RH}

\begin{theorem}[spectral census]\label{thm:census}
Let $n \ge m$. From the spectral data of $H_n(W)$ and the value $\mu_0(W)$:
\[
  \mu_0 = \text{number of zeros in } W \text{ with multiplicity},\qquad
  \rank H_n = m = r+2q,
\]
\[
  n_-(H_n) = q,\qquad n_+(H_n)-n_-(H_n) = r,
\]
and $\mu_0 > \sum_i 1 + \sum_j 2$ (i.e.\ $\mu_0$ exceeds the distinct count) if and only
if some zero in $W$ is multiple. Moreover, if $Q\in\R^{n\times m}$ has orthonormal
columns spanning $\ran H_n$, then the generalized eigenvalues of the compressed pencil
\[
  \widehat H_1:=Q^TH_n^{(1)}Q,\qquad \widehat H_0:=Q^TH_nQ
\]
are exactly the support points $x_i$ and $z_j,\bar z_j$.
\end{theorem}

\begin{proof}
The first four items are Theorem \ref{thm:inertia} together with
$\mu_0=\sum a_i+2\sum b_j$. For the pencil statement, enumerate the complex support as
$s_1,\ldots,s_m$ and its positive multiplicities as $d_1,\ldots,d_m$. Put
\[
 V=[v(s_1)\ \cdots\ v(s_m)],\qquad
 D=\operatorname{diag}(d_1,\ldots,d_m),\qquad
 \Sigma=\operatorname{diag}(s_1,\ldots,s_m).
\]
Then, with ordinary transpose (not conjugate transpose),
\[
 H_n=VDV^T,\qquad H_n^{(1)}=VD\Sigma V^T.
\]
The Vandermonde argument in Theorem \ref{thm:inertia} gives $\rank V=m$ and
$\ran_{\C}H_n=\ran_{\C}V$. Thus $B:=Q^TV$ is invertible and
\[
 \det(\widehat H_1-\lambda\widehat H_0)
 =\det(B)^2\det(D)\prod_{a=1}^m(s_a-\lambda).
\]
The compressed determinant is not identically zero and its roots, with multiplicity,
are precisely the support points. Conjugate support pairs make the determinant real.
\end{proof}

\begin{proposition}[symmetric quotient chart; Stieltjes form]\label{prop:stieltjes}
Let $0\le a<b$ and let $W^\pm$ denote the union of the two reflected ordinate windows
$(a,b)$ and $(-b,-a)$, with boundary contours chosen away from zeros. Define
$\mu_k(W^\pm)$ by adding the two contour integrals and set
$\nu_k(W^\pm):=\mu_{2k}(W^\pm)$. Then the zero multiset in $W^\pm$ is invariant under
$z\mapsto-z$, and $\nu_k$ is the moment sequence of the pushforward counting measure
under $w=z^2$. Under this map a real zero maps to $w\ge0$, while every nonreal orbit
produces either a nonreal conjugate pair in the $w$-plane or a negative real support
point. Consequently, for $n$ at least the number of distinct $w$-support points,
\[
  \text{all zeros of } A \text{ in } W^\pm \text{ are real}
  \iff
  \bigl(\nu_{\alpha+\beta}\bigr) \succeq 0
  \ \text{ and }\
  \bigl(\nu_{\alpha+\beta+1}\bigr) \succeq 0 ,
\]
the finite Stieltjes condition for the pushforward measure. Hence the zero-location
statement is equivalent to positivity of both even-moment Hankel matrices on every
member of a symmetric paired tiling.
\end{proposition}

\begin{proof}
The use of $W^\pm$, rather than a single generic interval, ensures that restriction to
the window commutes with the even quotient. Pairing the terms at $z$ and $-z$ gives a
positive integer weight at $w=z^2$. For a finite real moment sequence, positivity of
the Hankel matrix and its shift is equivalent to support in $[0,\infty)$; the
Vandermonde inertia argument of Theorem \ref{thm:inertia} supplies the converse directly
at any order at least the distinct support count. Finally $z^2\in[0,\infty)$ holds
exactly when $z\in\R$.
\end{proof}

\begin{corollary}[sign rigidity]\label{cor:rigidity}
Let $t\mapsto A_t$ be a continuous deformation with zeros varying continuously and no
zero crossing $\partial W$. At every parameter for which $n$ dominates the distinct
support count,
\[
 q(t)=n_-(H_n(t)).
\]
Thus the conjugate-pair count changes exactly when the negative inertia changes, and
such a change can occur only at a zero eigenvalue of $H_n(t)$. A tangency at zero need
not change the count. In particular, if $H_n(t)\succeq0$ for the entire deformation,
then $q(t)=0$ throughout. Positivity at a single parameter does not by itself propagate;
propagation requires exclusion of a subsequent negative-inertia change.
\end{corollary}

\section{Unconditional statements}

The following hold with no hypothesis.

\begin{enumerate}
\item[(U1)] For every window $W$ and $n\ge m(W)$: $q(W)=n_-\bigl(H_n(W)\bigr)$; the
off-critical-line pair count of a window is the negative inertia of a matrix computed by
contour integration of $A'/A$.
\item[(U2)] The census of Theorem \ref{thm:census} holds as stated.
\item[(U3)] The zero-location statement is equivalent to the Stieltjes condition of
Proposition \ref{prop:stieltjes} on every member of a symmetric paired tiling.
\item[(U4)] For every window contained in the numerically verified range
($|t|\le 3\cdot 10^{12}$ by the rigorous interval computation of Platt and
Trudgian~\cite{PlattTrudgian2021}), $q(W)=0$ and hence
$H_n(W)\succeq 0$: the positivity hypothesis is an established fact on an initial
segment of the tiling.
\item[(U5)] With $Y>\tfrac12$ and the convention
$A(z)=\xi(\tfrac12+iz)$, the bottom edge of $\partial R_W$ lies in the half plane of
absolute convergence and the top edge maps there by the functional equation, giving a
decomposition $H_n = H^{\mathrm{pr}} + H^{\mathrm{arch}} + H^{\mathrm{side}}$ into blocks
computable from Dirichlet data, $\Gamma$-factors and the lateral edges respectively; in a
tiling, adjacent lateral contributions cancel pairwise.
\end{enumerate}

\section{The Euler-anchored pencil and its regulator}

Nothing in this section refers to zeros; the dependency order is
\[
  \text{Euler/FE bank}\;\to\;(G_1,G_0)\;\to\;\text{real spectrum}\;\to\;
  \text{contour transport}\;\to\;\text{residues}\;\to\;\text{zeros},
\]
and no step may be inverted.

\begin{proposition}[Euler anchor]\label{prop:anchor}
Fix $s_0>1$ and $N\ge1$, and set
\[
  G_\ell[j,k] \;:=\; \sum_{n\ge2} \Lam(n)\,n^{-s_0}\,(\log n)^{\,j+k+\ell},
  \qquad 0\le j,k<N,\quad \ell\in\{0,1\}.
\]
The series converge absolutely, and $G_0,G_1$ are the Hankel matrices of the positive
measure $\sum_n \Lam(n)n^{-s_0}\delta_{\log n}$ supported on infinitely many points of
$[\log 2,\infty)$. Hence
\[
  G_0 > 0,\qquad G_1 > 0,
\]
and $\spec(G_1,G_0)\subset[\log2,\infty)$ consists of the Gauss quadrature nodes of that
measure.
\end{proposition}

\begin{proof}
Absolute convergence for $s_0>1$ is standard. Positivity is the Hamburger/Stieltjes
criterion for a positive measure on $[\log2,\infty)$: for $c\in\R^N$ and
$P(x)=\sum_j c_jx^j$, $c^{T}G_0c=\int P^2\,d\sigma\ge0$ and
$c^{T}G_1c=\int xP^2\,d\sigma\ge0$. If $c\ne0$, the polynomial $P$ has only finitely
many roots, whereas the prime-power logarithms in the support are infinite; both
integrals are therefore strictly positive.
\end{proof}

\begin{theorem}[finite-dimensional reverb regulator]\label{thm:regulator}
Let $t\mapsto (G_0(t),G_1(t))$ be a smooth family with $G_0(t)>0$ on the active quotient,
and set $S(t):=G_0^{-1/2}G_1G_0^{-1/2}$, so $S=S^{*}$ and
$\spec S=\spec(G_1,G_0)$. Suppose the active dimension is $d$, the spectrum of
$S(t_0)$ is simple, with spectral projections $P_a$ and unit eigenvectors $v_a$
($1\le a\le d$), and the transport has exactly $d$ scalar control channels,
\[
  \dot S \;=\; F(u) \;=\; F_0 + \sum_{r=1}^{d} u_r F_r ,
  \qquad F_r = F_r^{*}.
\]
Define the square response matrix $M\in\R^{d\times d}$ by
$M_{ar}:=\langle v_a,F_rv_a\rangle$ and put
$b_a:=-\langle v_a,F_0v_a\rangle$.
If $\det M \ne 0$, then $u := M^{-1}b$ is the unique control with
$\Delta\bigl(F(u)\bigr):=\sum_a P_aF(u)P_a = 0$, and with
\[
  K := \sum_{a\ne b}\frac{P_aF(u)P_b}{\lambda_a-\lambda_b},
  \qquad K^{*}=-K,
\]
one has $\dot S = [S,K]$. Consequently $S(t)=U(t)S(t_0)U(t)^{*}$ with $\dot U=-KU$,
$U$ unitary, and
\[
  \spec S(t) \;=\; \spec S(t_0)\quad\text{exactly.}
\]
In particular the transported pencil remains Hermitian-definite and its generalized
spectrum remains real.
\end{theorem}

\begin{proof}
First-order perturbation theory gives $\dot\lambda_a=\langle v_a,F v_a\rangle$, so
$\Delta(F)$ is precisely the component of the flow able to move eigenvalues; the linear
system $Mu=b$ expresses $\Delta(F(u))=0$ and is uniquely solvable iff $\det M\ne0$.
For such $u$, using $\sum_aP_a=I$,
\[
  [S,K]=\sum_{a\ne b}(\lambda_a-\lambda_b)\frac{P_aF(u)P_b}{\lambda_a-\lambda_b}
  =\sum_{a\ne b}P_aF(u)P_b=F(u)-\Delta(F(u))=F(u)=\dot S .
\]
$K^{*}=-K$ since $F(u)^{*}=F(u)$ and the denominators are real and antisymmetric in
$(a,b)$. The solution of $\dot U=-KU$, $U(t_0)=I$, is unitary, and
$\frac{d}{dt}\bigl(U^{*}SU\bigr)=U^{*}\bigl(\dot S-[S,K]\bigr)U=0$.
\end{proof}

\begin{remark}[numerical verification]
For the Euler anchor of Proposition \ref{prop:anchor} with $s_0=1.5$, $N=4$, control
channels given by reweighting the prime channels $p\in\{2,3,5,7\}$ and raw flow
$F_0=\partial_{s_0}S$: the spectrum is simple $(0.9619,2.7211,5.6399,8.8009)$,
$\det M=-0.128$ with $\rank M=4$ and condition number $3.4\cdot10^{2}$; the solved
control annihilates all diagonal drifts to $10^{-12}$, and
$\|[S,K]-(F-\Delta(F))\|=1.3\cdot10^{-11}$ with $\|K+K^{T}\|=1.7\cdot10^{-13}$.
\end{remark}

\begin{lemma}[block regulator at a degenerate spectrum]\label{lem:block}
Let $S=S^{*}$ have distinct eigenvalues $\Lambda_\alpha$ with spectral projections
$P_\alpha$ of arbitrary finite multiplicities, and let the flow be controlled as in
Theorem \ref{thm:regulator}, $\dot S=F(u)=F_0+\sum_ru_rF_r$, $F_r=F_r^{*}$. With
\[
  \Delta(F):=\sum_\alpha P_\alpha FP_\alpha,
  \qquad
  K:=\sum_{\alpha\ne\beta}\frac{P_\alpha FP_\beta}{\Lambda_\alpha-\Lambda_\beta},
\]
one has $[S,K]=F-\Delta(F)$ and $K^{*}=-K$. To first order the group $\Lambda_\alpha$
splits by the spectrum of the compression $P_\alpha FP_\alpha$; hence the regulated
condition is the block equation $P_\alpha F(u)P_\alpha=0$ for all $\alpha$, a linear
system in $u$ generalizing $Mu=b$. Whenever that block system is solvable on an interval
of fixed active dimension, the flow $\dot S=[S,K]$ is isospectral with multiplicities.
\end{lemma}

\begin{proof}
$SP_\alpha=\Lambda_\alpha P_\alpha$ and $\sum_\alpha P_\alpha=I$ give
\[
  [S,K]=\sum_{\alpha\ne\beta}(\Lambda_\alpha-\Lambda_\beta)
  \frac{P_\alpha FP_\beta}{\Lambda_\alpha-\Lambda_\beta}
  =F-\Delta(F),
\]
and $K^{*}=-K$ since $F^{*}=F$ and the real denominators are antisymmetric in
$(\alpha,\beta)$. The first-order motion of a degenerate group is degenerate
perturbation theory: the group splits by the eigenvalues of $P_\alpha FP_\alpha$,
which vanish identically iff the block equation holds; the condition is linear in
$u$, and on solution intervals the Lax argument of Theorem \ref{thm:regulator} is
unchanged.
\end{proof}

\begin{remark}[scope at singular events]
Lemma \ref{lem:block} handles eigenvalue multiplicity for a fixed-dimensional Hermitian
matrix. It does not by itself cross a rank drop of $G_0$, prove solvability of the block
response system, or extend the scalar interpolation warp through merging nodes. Those
three assertions require separate continuation data and are retained explicitly in the
terminal seating statement below.
\end{remark}

\section{Warp covariance: the bank-generated readout}\label{sec:warp}

Theorem \ref{thm:regulator} fixes the abstract spectrum $\{\lambda_a\}$ of the regulated
flow; the same transport, unregulated, moves it. The first chart makes reality manifest,
the second carries the physical ordinates, and the map between them --- the warp --- is
not a hypothesis to be assumed but an object the bank constructs. Throughout this
section $t\mapsto(G_0(t),G_1(t))$ is a given smooth transport of the continued jet pair, whose
inputs are values and derivatives of the continued logarithmic derivative along the path
--- function data; no zero location enters --- and the parameter interval is
\emph{regular}: $G_0(t)>0$ on the active quotient and $\spec S_t$ simple, where
$S_t:=G_0^{-1/2}G_1G_0^{-1/2}$. Theorems \ref{thm:warp} and \ref{thm:pullback}
apply on such regular intervals. Their hypotheses do not assert existence of the full
path, positivity of $G_0$ along it, or gluing through collision and rank-drop events.

\begin{definition}[bank velocity field; warp]\label{def:warp}
Let $\lambda_1(t)<\dots<\lambda_n(t)$ be the eigenvalues of $S_t$, realized as the
generalized eigenpairs $G_1w_a=\lambda_aG_0w_a$ normalized by $w_a^{T}G_0w_a=1$, and set
\[
  d_a(t) := w_a^{T}\bigl(\dot G_1-\lambda_a\dot G_0\bigr)w_a
  \qquad\bigl(=\langle v_a,\dot S_t\,v_a\rangle,\ v_a=G_0^{1/2}w_a\bigr).
\]
Let $V_t$ be the unique real polynomial of degree $\le n-1$ with
$V_t(\lambda_a(t))=d_a(t)$, and let the \emph{warp} $\varphi_t$ be the flow of the
scalar equation $\dot x=V_t(x)$, $\varphi_0=\mathrm{id}$.
\end{definition}

\begin{theorem}[warp existence and exact covariance]\label{thm:warp}
On every regular interval the warp is defined on a neighborhood of
$[\lambda_1(0),\lambda_n(0)]$ and satisfies:
\begin{enumerate}
\item[(i)] \textbf{bank determination:} $V_t$, hence $\varphi_t$, is computed from
$(G_0,G_1)(t)$ and $(\dot G_0,\dot G_1)(t)$ alone;
\item[(ii)] \textbf{reality and order:} $\varphi_t(\R\cap\mathrm{dom})\subseteq\R$ and
$x\mapsto\varphi_t(x)$ is strictly increasing, in particular injective;
\item[(iii)] \textbf{exact covariance:} $\varphi_t(\lambda_a(0))=\lambda_a(t)$ for
every $a$ and every $t$ in the interval.
\end{enumerate}
\end{theorem}

\begin{proof}
(i) is the construction. For (ii): $V_t$ has real coefficients --- real nodes, real
diagonal drifts of a Hermitian family --- so the flow preserves $\R$; the field is
polynomial in $x$ with coefficients continuous in $t$, hence locally Lipschitz in $x$
uniformly on compact $t$-intervals, so distinct scalar characteristics never cross and
$\varphi_t$ is strictly increasing where defined. For (iii): first-order perturbation
theory for the Hermitian-definite pencil gives
$\dot\lambda_a=w_a^{T}(\dot G_1-\lambda_a\dot G_0)w_a=d_a(t)=V_t(\lambda_a(t))$, so the
eigenvalue path $t\mapsto\lambda_a(t)$ solves the initial-value problem defining
$\varphi_t(\lambda_a(0))$, and uniqueness makes them equal. Existence: the node
characteristics are the globally defined eigenvalue paths; a point between consecutive
nodes is trapped between two of them by (ii); and continuous dependence on initial
conditions extends the flow to a neighborhood of the compact hull for the same times.
\end{proof}

\begin{theorem}[determinant pullback]\label{thm:pullback}
Let $p_t(x):=\det(S_t-xI)$, and let $I$ be an open real interval containing the
time-$0$ active spectrum on which the warp is defined. Then on $\varphi_t(I)$,
\[
  p_t \;=\; J_t\cdot\bigl(p_0\circ\varphi_t^{-1}\bigr),
\]
with $J_t$ continuous and nonvanishing; consequently
\[
  Z(p_t)\cap\varphi_t(I)\;=\;\varphi_t\bigl(\spec S_0\bigr)\;\subset\;\R .
\]
\end{theorem}

\begin{proof}
The field is polynomial in $x$, so $\varphi_t$ is a strictly increasing $C^1$
diffeomorphism of $I$ onto $\varphi_t(I)$. On $\varphi_t(I)$ the functions $p_t$ and
$p_0\circ\varphi_t^{-1}$ have the same zero set
$\{\varphi_t(\lambda_a(0))\}_a=\{\lambda_a(t)\}_a$ (covariance and injectivity,
Theorem \ref{thm:warp}), each zero simple (simple spectrum on a regular interval). The
ratio $J_t:=p_t/(p_0\circ\varphi_t^{-1})$ is continuous and nonvanishing off the common
zeros and extends across each with the nonzero limit
$p_t'(\lambda_a(t))\big/\bigl(p_0\circ\varphi_t^{-1}\bigr)'(\lambda_a(t))$.
\end{proof}

\begin{remark}[the two spectral charts]\label{rem:architectures}
The regulated and unregulated flows describe one transport in two charts: labels frozen
with a moving readout, or a moving spectrum read through the identity; by Theorem
\ref{thm:warp} their relation is $\varphi_t$. A constant-factor identity written with
one spectral variable cannot express this change of chart unless the two root multisets
already coincide. Moreover, when $n>m$, the full window determinant
$\det(H_n^{(1)}-\lambda H_n)$ vanishes identically because of the structural nullspace.
The correct finite determinant is the compressed determinant of Theorem
\ref{thm:census}; on every regular path its root motion is expressed by the pullback
identity of Theorem \ref{thm:pullback}.
\end{remark}

\begin{remark}[numerical verification]
On the convergent segment ($s_0\colon1.5\to1.2$, $N=4$, exact jets of $-\zeta'/\zeta$;
\texttt{tmp/att201\_warp\_probe.py}): the perturbation drift matches finite differences
of the eigenvalue paths to $8.9\cdot10^{-8}$. The interpolated flow carries
\[
 (1.1916,\,4.0395,\,9.6214,\,19.3384)
 \longmapsto
 (2.1714,\,9.2832,\,23.2371,\,47.5293)
\]
with maximal error $1.9\cdot10^{-10}$; an
interior test point remains strictly interior. The exact-jet anchor nodes differ from
those quoted after Theorem \ref{thm:regulator}, which realize the truncated prime sum at
finite cutoff; the two banks are distinct legitimate realizations, and the drift between
them is the truncation observation of the ledger.
\end{remark}

\section{The conditional theorem}

\begin{hypothesis}[$\mathrm{PSD}$]\label{hyp:psd}
For every window $W$ of a tiling of $\R$ and
$n=\max\{1,\mu_0(W)\}$, $H_n(W)\succeq0$. Equivalently, on every symmetric paired
window, both Hankel matrices of the even moments are positive semidefinite in the
quotient chart of Proposition \ref{prop:stieltjes}.
\end{hypothesis}

\begin{theorem}\label{thm:main}
Assume Hypothesis \ref{hyp:psd}. Then
\begin{enumerate}
\item every nontrivial zero of $\zeta$ lies on the critical line;
\item for each window the spectral data of $H_n(W)$ returns the full census of Theorem
\ref{thm:census}, with $n_-\equiv0$;
\item the zeros are the real generalized eigenvalues of the compressed, explicitly
computed symmetric pencils $(\widehat H_1(W),\widehat H_0(W))$ of Theorem
\ref{thm:census}.
\end{enumerate}
\end{theorem}

\begin{proof}
By Theorem \ref{thm:inertia}, $H_n(W)\succeq0$ with $n\ge m(W)$ forces $q(W)=0$: no
conjugate pair of zeros of $A$ lies in $W$. Nonreal zeros occur in conjugate pairs with
equal real part, hence both members lie in the same window of the tiling (for a pair on a
boundary, apply the argument to a shifted tiling). Every nontrivial zero of $\zeta$
corresponds to a zero of $A$ lying in some window, and only finitely many per window by
Riemann--von Mangoldt. Hence all zeros of $A$ are real, i.e.\ all nontrivial zeros of
$\zeta$ have real part $\tfrac12$. Items (2) and (3) are Theorem \ref{thm:census}; under
Hypothesis \ref{hyp:psd}, $\widehat H_0>0$ and $\widehat H_1$ is real symmetric.
\end{proof}

\section{The seat: endpoint identification}\label{sec:seat}

\begin{definition}[compressed terminal defect]\label{def:defect}
Fix a window $W$, let $m=\rank H_n(W)$, and choose
$Q_W\in\R^{n\times m}$ with orthonormal columns spanning $\ran H_n(W)$. Write
\[
 \widehat H_0:=Q_W^TH_n(W)Q_W,
 \qquad
 \widehat H_1:=Q_W^TH_n^{(1)}(W)Q_W.
\]
For a transported bank pencil whose terminal active dimension is $m$, identify its
terminal active coordinates with those of $Q_W$ and define
\[
 D_0:=G_0(1)-\widehat H_0,
 \qquad
 D_1:=G_1(1)-\widehat H_1.
\]
Both terms in each difference are obtained from function data: continued logarithmic-
derivative jets on the bank side and contour integrals of $A'/A$ on the window side.
No full $n\times n$ determinant is used; when $n>m$ that determinant contains a common
structural nullspace and is identically zero as a polynomial in the spectral parameter.
\end{definition}

\begin{hypothesis}[terminal seating $(\mathrm{S})$]\label{hyp:seat}
For every window $W$ there is a bank-generated transport from an Euler anchor to the
window chart, decomposed into finitely many regular intervals, with the following
properties:
\begin{enumerate}
\item on each regular interval $G_0>0$ on the active quotient and Theorem
\ref{thm:warp} supplies the real, order-preserving warp;
\item at every collision or rank change, the adjacent active spectra and warps admit
compatible continuous multiset limits;
\item the terminal multiset of those limits equals
$\spec(\widehat H_1,\widehat H_0)$.
\end{enumerate}
Every object used to construct the path and the transition data is required to come
from Euler, functional-equation, jet, contour, and registration data, without inserting
the support points as interpolation targets.
\end{hypothesis}

\begin{theorem}[the conditional chain]\label{thm:chain}
Hypothesis \ref{hyp:seat} implies Hypothesis \ref{hyp:psd}, and therefore implies all
three conclusions of Theorem \ref{thm:main}.
\end{theorem}

\begin{proof}
On each regular interval the generalized eigenvalues are real because $G_0>0$ and
$G_1$ is real symmetric; equivalently, Theorem \ref{thm:warp} carries the real anchor
spectrum by a real map. The transition clause of Hypothesis \ref{hyp:seat} preserves
reality under multiset limits. Its terminal clause therefore makes every root of the
compressed determinant
$\det(\widehat H_1-\lambda\widehat H_0)$ real. By Theorem \ref{thm:census}, these roots
are exactly the zero support in $W$, so $q(W)=0$. Theorem \ref{thm:inertia} now gives
$H_n(W)\succeq0$ for every window, which is Hypothesis \ref{hyp:psd}.
\end{proof}

\begin{remark}[what the warp theorem proves]
Theorems \ref{thm:warp} and \ref{thm:pullback} prove an exact covariance statement for
every specified regular Hermitian-definite path. They do not prove the existence of a
path reaching the contour-moment pencil, continued positivity of $G_0$, compatibility
at singular events, or terminal equality with the compressed window spectrum. Those
claims are precisely the content of Hypothesis \ref{hyp:seat}; coordinate covariance
alone is not endpoint identification.
\end{remark}

\section{Remaining analytic obligation}\label{sec:obligations}

The endpoint comparison is to be performed on the compressed active pencils. Its
computable entrywise form is the following defect evaluation:
\[
  G_\ell[j,k]
  =\frac{1}{2\pi i}\oint_\Gamma
    z^{\,j+k+\ell}\frac{A'(z)}{A(z)}\,dz
   +D_\ell[j,k],
  \qquad \ell\in\{0,1\},
\]
where $D_\ell$ must be derived explicitly from the pole, $\Gamma$-factor, lateral-edge,
windowing, and coordinate-compression terms. On $s_0>1$, the pole contribution has
moments $k!/(s_0-1)^{k+1}$, the moments of
$e^{-(s_0-1)x}\,dx$ on $[0,\infty)$, and is positive semidefinite. The signs and spectral
effect of the remaining blocks must be retained rather than absorbed into a constant
factor.

Once this identity and the event-gluing data establish Hypothesis \ref{hyp:seat},
Theorem \ref{thm:chain} applies. Until then, the unconditional results are the inertia
identity, the compressed spectral census, Euler-anchor positivity, the regulator under
its stated response hypotheses, and regular-interval warp covariance.

\begin{thebibliography}{9}
\bibitem{PlattTrudgian2021}
D. Platt and T. Trudgian,
\emph{The Riemann hypothesis is true up to $3\cdot10^{12}$},
Bulletin of the London Mathematical Society \textbf{53} (2021), 792--797,
\href{https://doi.org/10.1112/blms.12460}{doi:10.1112/blms.12460}.
\end{thebibliography}

\end{document}
